

% article class because we want to fully customize the page and not use a cv template
\documentclass[a4paper,11pt]{article}

\usepackage{url}
\usepackage{parskip} 	

%other packages for formatting
\RequirePackage{color}
\RequirePackage{graphicx}
\usepackage[usenames,dvipsnames]{xcolor}
% \usepackage[scale=0.9]{geometry}
\usepackage[margin=0.4in]{geometry}



%tabularx environment
\usepackage{tabularx}

%for lists within experience section
\usepackage{enumitem}

% centered version of 'X' col. type
\newcolumntype{C}{>{\centering\arraybackslash}X} 

%to prevent spillover of tabular into next pages
\usepackage{supertabular}
\usepackage{tabularx}
\newlength{\fullcollw}
\setlength{\fullcollw}{0.47\textwidth}

\usepackage{bibentry}
\makeatletter\let\saved@bibitem\@bibitem\makeatother
% \usepackage[colorlinks=true]{hyperref} % loaded later with consistent options
\makeatletter\let\@bibitem\saved@bibitem\makeatother


% hanging publications: 1st line starts at beginning of line, 
% further lines are placed a bit to the right
\usepackage{hanging}
\newcommand\publication[1]{%
	\smallskip\par\hangpara{1.5em}{1}\bibentry{#1}\smallskip
}


%custom \section
\usepackage{titlesec}				
\usepackage{multicol}
\usepackage{multirow}
\usepackage{xcolor}
\usepackage[dvipsnames]{xcolor}
\usepackage{parskip} 
\usepackage{relsize}
\usepackage{scrextend}

%CV Sections inspired by: 
%http://stefano.italians.nl/archives/26
\titleformat{\section}{\large\scshape\raggedright}{}{0em}{}[\titlerule]
\titlespacing{\section}{0pt}{6.5pt}{6.5pt}

%for publications
\usepackage[style=authoryear,sorting=ynt, maxbibnames=2]{biblatex}

%Setup hyperref package, and colours for links
\usepackage[unicode, draft=false]{hyperref}
\definecolor{linkcolour}{rgb}{0,0.2,0.6}
\hypersetup{colorlinks,breaklinks,urlcolor=linkcolour,linkcolor=linkcolour}
\addbibresource{citations.bib}
\setlength\bibitemsep{1em}

%for social icons
\usepackage{fontawesome5}

\usepackage{lmodern}



%debug page outer frames
%\usepackage{showframe}

%----------------------------------------------------------------------------------------
%	BEGIN DOCUMENT
%----------------------------------------------------------------------------------------
\begin{document}

% non-numbered pages
\pagestyle{empty} 



\begin{tabularx}{\linewidth}{@{} C @{}}
  \Huge {\textbf{Obed Junias} }\\[1pt]
  \href{mailto:obed.junias@colorado.edu}{\raisebox{-0.05\height}\faEnvelope \ obed.junias@colorado.edu} \ $|$ \
  \href{https://obedjunias.github.io}{\raisebox{-0.05\height}\faGlobe\ obedjunias.com} \ $|$ \
  % \href{https://github.com/obedjunias19?tab=repositories}{\raisebox{-0.05\height}\faGithub \ obedjunias19} \ $|$ \
  \href{https://github.com/obedjunias19}{\raisebox{-0.05\height}\faGithub \ obedjunias19} \ $|$ \
  \href{https://linkedin.com/in/obed-junias}{\raisebox{-0.05\height}\faLinkedin\ obed-junias} \ $|$ \
  \href{tel:+1720-266-0046}{\raisebox{-0.05\height}\faMobile \ 720-266-0046} \\
\end{tabularx}


\color{BlueViolet}
\section*{Research Summary}
\color{black}
I study \textbf{model collapse in language models}, focusing on how recursive training on model-generated data degrades distributional diversity, factual reliability, and long-term model quality. My current work combines \textbf{measurement, detection, and controlled experimentation} to identify early signals of collapse under synthetic-data feedback loops.

More broadly, I am interested in reliable reasoning, evaluation under distribution shift, and scalable systems for studying foundation models.

\textit{Long-term goal: developing principled methods to measure, detect, and mitigate model collapse in foundation models across text and code generation.}

\color{BlueViolet}
\section*{Primary Research Focus}
\color{black}
Model Collapse in Generative AI, Synthetic Data, Recursive Training, Reliability of Foundation Models

\color{BlueViolet}
\section*{Secondary Interests}
\color{black}
Commonsense Reasoning, LLM Evaluation and Benchmarking, Responsible AI, Scalable LLM Systems


%----------------------------------------------------------------------------------------
%	EDUCATION
%----------------------------------------------------------------------------------------
\color{BlueViolet}
\section*{Education}
\color{black}
\begin{tabularx}{\linewidth}{@{}X r@{}}
\textbf{M.S. in Computer Science}, \textbf{University of Colorado Boulder} & Aug 2024 – May 2026 \\
\textit{Courses:} NLP, Deep Natural Language Understanding, Neurosymbolic NLP Methods & (GPA: 4.0/4.0)
\end{tabularx}
\begin{tabularx}{\linewidth}{@{}X r@{}}
\textbf{B.E. in Computer Science}, BMS College of Engineering & Aug 2017 – Aug 2021 \\
\textit{Relevant Courses:} Algorithms, Databases, Operating Systems, Machine Learning
\end{tabularx}

%----------------------------------------------------------------------------------------
%\tPUBLICATIONS
%----------------------------------------------------------------------------------------
\color{BlueViolet}
\section{Publications}
\color{black}
\begin{tabularx}{\linewidth}{ @{}l r@{} }
\multicolumn{2}{@{}X@{}}{
\begin{minipage}[t]{\linewidth}
\begin{itemize}[nosep, after=\strut, leftmargin=1.2em, itemsep=1.5pt]
\item \textbf{Junias, Obed}, and Maria Leonor Pacheco. \textit{LOGICAL-COMMONSENSEQA: A Benchmark for Logical Commonsense Reasoning.} \textbf{ACL 2026 (Main Conference)} \href{https://arxiv.org/abs/2601.16504}{[paper]}

\item \textbf{Junias, Obed*}, et al. \textit{Assessing Algorithmic Bias in Language-Based Depression Detection: A Comparison of DNN and LLM Approaches.} \textbf{IEEE-EMBS International Conference on Biomedical and Health Informatics (IEEE-BHI), 2025} \href{https://ieeexplore.ieee.org/abstract/document/11269509}{[paper]}
\end{itemize}
\end{minipage}
}
\end{tabularx}

%----------------------------------------------------------------------------------------
% EXPERIENCE SECTIONS
%----------------------------------------------------------------------------------------

%Industry Experience
\color{BlueViolet}
\section{Research Experience}
\color{black}

\begin{tabularx}{\linewidth}{@{}X r@{}}
\textbf{Measuring Model Collapse under Recursive Summarization Training} & Feb 2026 - Present \\
\multicolumn{2}{@{}X@{}}{\textit{Advisor: \href{https://blast-cu.github.io/mlpacheco/}{Maria Leonor Pacheco}, \href{https://blast-cu.github.io}{BLAST Lab}}} \\
\end{tabularx}
\begin{itemize}[nosep, after=\strut, leftmargin=1.5em, itemsep=2pt]
    \item Designing a controlled recursive training pipeline where human-written corpora are progressively replaced with \textbf{model-generated (synthetic) data} across generations.
    \item Training language models under synthetic-data feedback loops to study degradation in \textbf{distributional diversity, factual reliability, and knowledge fidelity}.
    \item Developing measurement methods for collapse using \textbf{entropy-based metrics, lexical diversity, and distributional divergence}.
    \item Investigating early warning signals such as \textbf{entropy decay, repetition, and detectability of synthetic text} as indicators of collapse before major downstream failures.
\end{itemize}

\begin{tabularx}{\linewidth}{@{}X r@{}}
\textbf{Structured Commonsense Reasoning and Benchmarking} & June 2025 – Present \\
\multicolumn{2}{@{}X@{}}{\textit{Advisor: \href{https://blast-cu.github.io/mlpacheco/}{Maria Leonor Pacheco}, \href{https://blast-cu.github.io}{BLAST Lab}, Associated publication listed above}} \\
\end{tabularx}
\begin{itemize}[nosep, after=\strut, leftmargin=1.5em, itemsep=2pt]
    \item Introduced \textbf{LOGICAL-COMMONSENSEQA}, a benchmark designed to isolate compositional and multi-fact reasoning failures in LLMs.
    \item Identified systematic mismatches between \textbf{chain-of-thought fluency and logical correctness}, demonstrating cases where coherent reasoning traces yield incorrect conclusions.
    \item Developed structured evaluation protocols to analyze reasoning depth and failure modes.
    \item Designed entailment-like reasoning representations to study how performance varies with task structure and inference depth.
\end{itemize}


\begin{tabularx}{\linewidth}{@{}X r@{}}
\textbf{Responsible LLM Evaluation in Mental Health Applications} & Jan 2025 – Nov 2025 \\
\multicolumn{2}{@{}X@{}}{\textit{Advisor: \href{https://www.colorado.edu/cs/theodora-chaspari}{Theodora Chaspari}, HUBBS Lab, Associated publication listed above}} \\
\end{tabularx}
\begin{itemize}[nosep, after=\strut, leftmargin=1.5em, itemsep=2pt]
    \item Built \textbf{large-scale LLM evaluation pipelines} to assess demographic bias and subgroup performance of LLMs (LLaMA, GPT-4) for mental health classification tasks.
    \item Studied how prompting strategies and fairness-aware losses affect reasoning consistency and robustness in low-resource and imbalanced data settings.
    \item Developed diagnostic metrics to characterize bias and reasoning instability beyond aggregate accuracy.
    \item Led research design for a \textbf{peer-reviewed publication}, demonstrating limitations of LLM reliability in sensitive domains.
\end{itemize}


\begin{tabularx}{\linewidth}{@{}X r@{}}
\textbf{Efficient Multi-Agent LLM Inference with Parameter-Efficient Adaptation} & Aug 2025 – Dec 2025 \\
\end{tabularx}
\begin{itemize}[nosep, after=\strut, leftmargin=1.5em, itemsep=2pt]
    \item Designed a modular inference framework to study how parameter-efficient adapters enable scalable multi-agent LLM systems.
    \item Implemented dynamic \textbf{LoRA adapter loading and scheduling} under strict GPU memory constraints, streaming adapters from CPU to GPU at runtime.
    \item Conducted controlled experiments comparing adapter-based specialization against full-model switching, analyzing latency, throughput, and memory efficiency.
    \item Demonstrated that adapter-based agents reduce cumulative switching overhead by \textbf{20–90$\times$}, enabling complex agent workflows on constrained hardware.
\end{itemize}

\color{BlueViolet}
\section{Relevant Research Projects}
\color{black}
\begin{tabularx}{\linewidth}{@{}X r@{}}
\textbf{Lost in Plot: Contrastive Learning for Logical Semantic Retrieval} & \\
\end{tabularx}
\begin{itemize}[nosep, after=\strut, leftmargin=1.5em, itemsep=2pt]
    \item Formulated a semantic retrieval task for vague queries, targeting underspecified and ambiguous natural-language descriptions requiring semantic rather than lexical matching.
    \item Designed a dual-encoder model (BERT) with multi-task contrastive learning objectives to align free-form user descriptions with structured movie metadata.
    \item Constructed a 100K+ synthetic dataset to enable scalable training and controlled evaluation.
    \item Evaluated retrieval quality using Recall@K and MRR, outperforming a GPT-4 few-shot baseline.
\end{itemize}


\begin{tabularx}{\linewidth}{@{}X r@{}}
\textbf{Medical Ethics Assessment of Large Language Models} & \\
\end{tabularx}
\begin{itemize}[nosep, after=\strut, leftmargin=1.5em, itemsep=2pt]
    \item Designed evaluation protocols to assess ethical reasoning of LLMs in medical decision-making scenarios.
    \item Built a multiple-choice benchmarking pipeline to analyze how external knowledge affects reasoning consistency.
    \item Analyzed rationales and confidence signals to identify systematic failure modes in ethically sensitive settings.
\end{itemize}


\color{BlueViolet}
\section{Industry Experience}
\color{black}
\begin{tabularx}{\linewidth}{ @{}l r@{} }
\textbf{Senior Member of Technical Staff} at  Oracle Corporation & \hfill Aug 2021 - Aug 2024 \\[3.75pt]
\multicolumn{2}{@{}X@{}}{
\begin{minipage}[t]{\linewidth}
    \begin{itemize}[nosep,after=\strut, leftmargin=1.5em, itemsep=2pt]
        \item Designed and implemented scalable automation and validation frameworks (\textbf{TestNG, BATS}) for end-to-end testing of \href{https://www.oracle.com/database/technologies/appdev/rest.html} {Oracle REST Data Services (ORDS)} and \href{https://www.oracle.com/database/sqldeveloper/technologies/sqlcl/}{SQLcl}, improving validation efficiency by \textbf{~95\%}.       
        \item Led technical development, validation and delivery of a new Oracle Cloud \href{https://docs.oracle.com/en/cloud/paas/autonomous-database/serverless/adbsb/autonomous-tools.html#GUID-8B5AAB5D-21B7-4EFC-8018-2AE18642FAFB}{DB Tools} feature, coordinating a 15-engineer effort and achieving production readiness within 2 months.
    \end{itemize}
    \end{minipage}
}
\end{tabularx}


\begin{tabularx}{\linewidth}{ @{}l r@{} }
\textbf{Machine Learning Engineer Intern} at  Hewlett Packard Enterprise & \hfill Feb 2021 - Jul 2021 \\ 
%\textit{Tech Stack: Python, Groovy, Workfusion} \\[3.75pt]
\multicolumn{2}{@{}X@{}}{
\begin{minipage}[t]{\linewidth}
    \begin{itemize}[nosep,after=\strut, leftmargin=1.5em, itemsep=2pt]
        \item Built OCR-based automation pipelines for processing unstructured enterprise documents, integrating learned document extraction into production workflows.
        \item Evaluated model performance across noisy and heterogeneous inputs, contributing to improvements in robustness and downstream automation reliability.
    \end{itemize}
    \end{minipage}
}
\end{tabularx}

\color{BlueViolet}
\section{Skills}
\color{black}

\begin{tabularx}{\linewidth}{@{}l X@{}}
  \textbf{Languages \& Frameworks} & Python, PyTorch, Triton, LangGraph, SQL, Transformers, C++\\
  % \textbf{ML Techniques} & Regression, Naive Bayes, Random Forests, SVMs, Neural Networks \\
  \textbf{LLM Systems} & LLM Serving, LLM Evaluation, RAG Pipelines, Multi-Agent Systems and Tools \\
  \textbf{LLM Training} & Pre-Training, Supervised Fine-Tuning, Transformers, PyTorch DDP \\
  \textbf{ML Infrastructure} & CUDA, vLLM, FlashAttention, Quantization, GPU Memory Optimization \\


\end{tabularx}

\color{BlueViolet}
\section*{Awards and Recognition}
\color{black}
\begin{tabularx}{\linewidth}{ @{}l r@{} }
\multicolumn{2}{@{}X@{}}{
\begin{minipage}[t]{\linewidth}
\begin{itemize}[nosep, after=\strut, leftmargin=1.2em, itemsep=1.5pt]
    \item \textbf{NSF--EMBS--Google Young Professional NextGen Scholar}, \textit{IEEE EMBS BHI Conference}, 2025
\end{itemize}
\end{minipage}
}
\end{tabularx}



\vfill

\end{document}
